\scanpage{232}
et de la structure de graphe $(S\deleted{\unsure{_0}},A)$, et non
de la famille $\pi$ de sous-ens $A_i\subset A$, $i\in\pi$.
Ceci posé, on a

\deleted{(44) Les cond. équivalentes (41) (41') (41'') signifient
aussi que pour deux \inserted{$Z$-}rayons incidents,
$D,D'\subset S$, et pour tt $(s,s')\in D\times D'$, $\exists$ existe
un \uline{unique} $\deleted{\unsure{\lambda}}i=i_{s,s'}\in\pi$ tel que $\{s,s'\}\in A_i$, \deleted{On} \inserted{soit}
\deleted{de \unsure{int},} \deleted{a alors} On a $i_{\lambda s,s'}=i_{s,\lambda s'}$, et on trouve
ainsi un iso. canonique}
\[
 \deleted{D\wedge_ZD'\overset{\tau}{\simeq}\pi}\tag*{\deleted{(45)}}
\]
\deleted{\unsure{caractérisé|caractérisée} par $i_{s,s'}=\tau(s\wedge_Zs')$, i.e. par la \unsure{condition|égalité|relation}}
\[
 \deleted{\unsure{i_{s,s'}}}\tag*{\deleted{(46)}}
\]
(44) Les conditions \deleted{(41)} équivalentes (41), (41') (41'')
signifient aussi que les $A_i$ sont mutuellement
disjoints, i.e. \deleted{\unsure{on}} \inserted{on} a une application
canonique ``type''
% slanted
\margin{donc la structure de multigraphes $(A_i)_{i\in\pi}$ sur $S$ (avec un groupe d'homothéties $Z$)
est connue quand on connaît $S,A\deleted{\unsure{\bigcup_{i\in\pi}}}\subset\mathcal P_{\{1,2\}}(S)$,
sa partition en $A=\coprod_{i\in\pi}A_i$, et enfin le sous-groupe $Z\subset\unsure{\mathfrak S|\mathfrak G}_S$,
(avec la condition que $Z$ est formé d'homothéties pour les structures de graphe $(S,A_i)$, et que
\unsure{l'application} $\unsure{Z\ldots}$ transforme \unsure{canoniq.} les unes \unsure{dans les autres})}
\deleted{On trouve alors}
\[
 A\xrightarrow{\tau_A}\pi
 \tag*{\deleted{(45)}}
\]
\deleted{Si cette la} Si \inserted{$Z$} $D,D'$ sont \deleted{deux} rayons incidents,
\[
 (s,s')\longmapsto\tau_A(\{s,s'\})\qquad D\times D'\xrightarrow{\tau_{D,D'}}\pi
\]
satisfait
\[
 \deleted{\tau_A(\{\lambda s,s'\})}\quad
 \tau_{D,D'}(\lambda s,\mu s')=\deleted{\tau_{D,D'}(s,\lambda s')}
 \inserted{\lambda\mu\,\tau_{D,D'}(s,s')}, \tag*{(46)}
\]
et induit \deleted{donc} un isom. de $Z$-torseurs
\[
 D\wedge D'\xrightarrow{\sim}\pi .\tag*{(47)}
\]
% margin note written sideways
\margin{\deleted{Dans le cas}}
\uline{NB} La donnée d'un prégraphe $\Gamma$, avec groupe
$Z$ commutatif d'automorphismes -- jouant le
rôle d'homothéties virtuelles -- permet de
définir dans $S$ (disons, \unsure{et il suffit|et il faut|bien sûr} \deleted{\unsure{pour|plus}} $S^\vee$)
% archive sheet number: 231

